\documentclass[11pt,a4paper]{article}
\usepackage[utf8]{inputenc}
\usepackage[T1]{fontenc}
\usepackage[margin=2.2cm]{geometry}
\usepackage{tabularx}
\usepackage{booktabs}
\usepackage{xcolor}
\usepackage{colortbl}
\usepackage{hyperref}
\usepackage{enumitem}
\usepackage{tcolorbox}
\usepackage{fancyhdr}
\hypersetup{colorlinks=true,linkcolor=black,urlcolor=blue!60!black}

\definecolor{rowgray}{RGB}{243,245,248}
\definecolor{good}{RGB}{18,105,58}
\definecolor{bad}{RGB}{158,32,32}
\definecolor{warn}{RGB}{176,110,10}
\definecolor{navy}{RGB}{28,42,64}

\newcolumntype{Y}{>{\raggedright\arraybackslash}X}
\newcommand{\SHOW}{\textcolor{good}{\textbf{SHOW}}}
\newcommand{\HIDE}{\textcolor{bad}{\textbf{DO NOT SHOW}}}
\newcommand{\CARE}{\textcolor{warn}{\textbf{SHOW WITH CARE}}}

\newtcolorbox{keybox}[1]{colback=blue!3,colframe=navy,fonttitle=\bfseries,title={#1}}
\newtcolorbox{warnbox}[1]{colback=orange!4,colframe=warn,fonttitle=\bfseries,title={#1}}

\pagestyle{fancy}\fancyhf{}
\fancyhead[L]{\small RAFT --- what it is, why we chose it, what we measured}
\fancyfoot[C]{\thepage}

\title{\vspace{-1.6cm}\textbf{RAFT for a Graph-RAG Assistant}\\[0.25cm]
\large Everything behind the fine-tuning decision:\\the method, the steps, the metrics, and what is safe to present}
\author{Prepared for Iyed Mediouni --- PFE, TEK-UP $\times$ TELNET Holding}
\date{17 August 2026}

\begin{document}
\maketitle
\thispagestyle{fancy}

\begin{keybox}{How to read this document}
Every number in here was measured on your system. Each result carries a flag:
\SHOW{} means it is solid and defensible under questioning; \CARE{} means it is real
but needs a stated caveat; \HIDE{} means it would not survive a sharp question and
should stay out of the report and the defence.

Sections 2--4 teach the method. Sections 5--7 are what we did and what came out.
Section 8 is the honest limitations list. Section 9 is defence preparation.
\end{keybox}

\tableofcontents
\newpage

\section{The problem in one page}

Your assistant answers questions about internal documentation. Retrieval finds
passages; the language model writes an answer from them. Two things can go wrong,
and they are completely different failures:

\begin{enumerate}[leftmargin=1.4em]
  \item \textbf{Retrieval does not find the answer.} No amount of model training
        fixes this --- the model never sees the text.
  \item \textbf{Retrieval finds it and the model does not use it.} No amount of
        code fixes this --- reading is what the model does.
\end{enumerate}

The whole fine-tuning argument rests on separating those two. If you cannot
separate them, any claim about fine-tuning is contestable, because a reviewer can
always answer ``you should have fixed retrieval instead.''

We measured both.

\begin{table}[h]\centering\small
\begin{tabularx}{\textwidth}{Y r l}
\toprule
\textbf{Measurement (500-question held-out set)} & \textbf{Value} & \textbf{Whose fault} \\
\midrule
Retrieval returned the answering passage & 16\,\% & retrieval \\
\rowcolor{rowgray} Model used it when it was ranked 1st--2nd & 71.6\,\% & --- \\
Model ignored it entirely when ranked 1st--2nd & \textbf{28.4\,\%} & \textbf{the model} \\
\rowcolor{rowgray} Answer contained a fabricated identifier & 2.6\,\% & the model (rare) \\
Answer quoted its source verbatim & 2.19\,\% & the model \\
\bottomrule
\end{tabularx}
\end{table}

The third row is the one that justifies fine-tuning. A passage ranked \emph{first}
cannot be blamed on ordering, cannot be recovered by raising \texttt{top\_k}, and
cannot be repaired by a post-processing filter. The model was shown the answer and
wrote something that contained none of it, more than a quarter of the time.

\section{What RAFT actually is}

\subsection{The exam analogy}

The RAFT authors frame it as exam preparation, and the analogy is worth keeping
because it makes the design obvious.

\begin{itemize}[leftmargin=1.4em]
  \item \textbf{Closed-book exam.} The model answers from memory alone. This is a
        plain language model with no retrieval.
  \item \textbf{Open-book exam.} The model is handed reference pages and answers
        from them. This is RAG --- your system.
  \item \textbf{Studying for an open-book exam.} Practising with the book open,
        including pages that look relevant but are not, until you get good at
        finding the right passage quickly. \textbf{This is RAFT.}
\end{itemize}

The key insight: standard RAG never practises. The model was trained to write
fluent text, then at inference time is suddenly handed six documents and expected
to pick correctly among them. RAFT trains that specific skill.

\subsection{The training recipe}

Each training example is a question, a set of documents, and an answer. Two kinds
of example, mixed in a fixed proportion $P$:

\begin{table}[h]\centering\small
\begin{tabularx}{\textwidth}{l Y}
\toprule
\textbf{$P$\,\% of examples} & question + \textbf{golden} passage + several \textbf{distractor} passages
$\rightarrow$ answer that reasons and quotes the golden passage \\
\rowcolor{rowgray} \textbf{$(1-P)$\,\% of examples} & question + \textbf{distractors only}, no golden passage
$\rightarrow$ (in our version) an explicit abstention \\
\bottomrule
\end{tabularx}
\end{table}

\textbf{Golden document:} the passage that actually answers the question.\\
\textbf{Distractor document:} a passage that is topically related but does not
contain the answer. In your corpus these are brutal --- ask about release 4.16.0
and retrieval returns 4.16.1, 4.20.3 and 4.28.0. They look almost identical.

We used $P = 0.80$: 794 examples with the golden passage present, 199 without.

\subsection{The chain-of-thought format}

RAFT does not train the model to emit a bare answer. It trains a reasoning chain
that quotes its evidence:

\begin{quote}\ttfamily\small
\#\#Reason: The documentation states that the 2.5.0 micro-service uses an ingress
controller. \#\#begin\_quote\#\# File Micro Service version 2.5.0 integrates a
HTTP reverse proxy and uses an ingress controller to expose only one entry point
\#\#end\_quote\#\# Therefore the ingress controller must be installed first.\\
\#\#Answer: Install the ingress controller before deploying, and keep previously
exposed services as ClusterIP rather than LoadBalancer.
\end{quote}

Two things happen when a model is trained on thousands of these:

\begin{enumerate}[leftmargin=1.4em]
  \item \textbf{It learns to locate before it writes.} The quote must come from
        somewhere, so the model develops the habit of finding the supporting span
        first.
  \item \textbf{The answer becomes checkable.} A reader can verify the quote
        against the page. An unquoted answer has to be trusted.
\end{enumerate}

The RAFT paper shows this matters: training on answers alone (no reasoning chain)
causes loss to fall quickly and the model to overfit to short answers. The
reasoning chain prevents that.

\subsection{Why the distractor-only examples exist}

This is the subtle part, and it is where we deliberately diverged from the paper.

In the original RAFT, the $(1-P)$ examples \emph{keep the same answer} even though
the golden document was removed. The paper's words: it ``compel[s] the model to
memorize answers instead of deriving them from the context.'' For their benchmarks
that is a feature --- it builds robustness.

For a documentation assistant it is a bug. A model taught to answer when the
evidence is absent is a model taught to fabricate. This weakness is documented:
\emph{Divide-Then-Align} (ACL 2025) shows RAFT ``conditions models to generate
answers even in the absence of reliable knowledge.''

\begin{keybox}{{Our deviation, and why it is defensible}}
In our training set the distractor-only examples target an \textbf{abstention} ---
``the provided context does not contain this information'' --- rather than the
memorised answer. The model is taught: \emph{quote the passage when it is there,
decline when it is not.}

This aligns with Google Research's ``sufficient context'' work, which recommends
treating abstention as a first-class response rather than a failure to minimise.
State this deviation explicitly in the report; it shows you read the method
critically rather than copying it.
\end{keybox}

\subsection{What the published results show}

\begin{table}[h]\centering\small
\begin{tabular}{l ccccc}
\toprule
 & \textbf{PubMed} & \textbf{HotpotQA} & \textbf{HuggingFace} & \textbf{Torch Hub} & \textbf{TensorFlow} \\
\midrule
LLaMA2-7B + RAG & 58.8 & 0.03 & 26.43 & 8.60 & 43.06 \\
\rowcolor{rowgray} Domain fine-tune (DSF) & 59.7 & 6.38 & 61.06 & 84.94 & 86.56 \\
DSF + RAG & 71.6 & 4.41 & 42.59 & 82.80 & 60.29 \\
\rowcolor{rowgray} \textbf{RAFT} & \textbf{73.30} & 35.28 & \textbf{74.00} & \textbf{84.95} & \textbf{86.86} \\
\bottomrule
\end{tabular}
\end{table}

Look at the third row. \textbf{Domain fine-tuning plus RAG is often worse than
domain fine-tuning alone} --- HuggingFace drops 61 to 43, TensorFlow 87 to 60.
Naively fine-tuning a model and then bolting retrieval onto it actively hurts.

That single row is your strongest published answer to ``why not just fix the
code'': the interaction between fine-tuning and retrieval is a real research
problem with a published remedy, not a plumbing detail.

\section{Why RAFT and not the alternatives}

\begin{table}[h]\centering\small
\begin{tabularx}{\textwidth}{l l Y}
\toprule
\textbf{Method} & \textbf{Year} & \textbf{Why it was not chosen} \\
\midrule
Vanilla QLoRA SFT & 2023 & Teaches answering style, not document selection. Does not
address the measured defect. \\
\rowcolor{rowgray} DPO on corrupted pairs & 2023 & We built this and the gate rejected the adapter.
The pairs were 93\,\% synthetic corruptions; the model learned to decline. \\
GRPO & 2024--25 & Needs a learned or online reward. Ours is already exact,
deterministic and free --- adding a reward model adds approximation error where an
exact reward exists. \\
\rowcolor{rowgray} RLHF / PPO & 2022 & Requires human preference labels nobody at TELNET produced. \\
ORPO / SimPO / KTO & 2024 & Preference variants; same data problem as DPO, and each
needs a new trainer wired into the pipeline. \\
\textbf{RAFT} & \textbf{2024} & \textbf{Targets document selection --- the measured
defect --- and runs on the existing QLoRA trainer with only the dataset changed.} \\
\bottomrule
\end{tabularx}
\end{table}

\begin{keybox}{The distinction to get right in the defence}
RAFT is \textbf{not} an alternative to QLoRA. QLoRA is \emph{how} the weights are
updated efficiently; RAFT is \emph{what} the model is shown. You ran RAFT
\emph{with} QLoRA. Anyone who knows the field will notice if you present them as
competing options.
\end{keybox}

\section{RAFT or GraphRAFT? --- and why traversal scored zero here}

This deserves its own section because the answer is counter-intuitive and a jury
may well ask.

\subsection{GraphRAFT does not apply to this architecture}

GraphRAFT (2025) extends RAFT to knowledge graphs. Its contribution is training a
model to \textbf{generate correct Cypher or SPARQL queries} for multi-hop graph
questions --- the LLM writes the query, the graph executes it.

\textbf{Your system does not work that way.} In this architecture,
\texttt{rag/neo4j\_search.py} performs the traversal itself, in Python, with
hand-written Cypher: vector search finds seed nodes, then the graph is expanded
along typed relationships (\texttt{BLOCKS}, \texttt{RELATES\_TO},
\texttt{CHILD\_OF}, \texttt{CAUSES}, \dots), then a cross-encoder reranks. The
language model never sees a query language; it receives finished passages.

Adopting GraphRAFT would therefore mean \emph{re-architecting retrieval} so that
the LLM generates the queries --- replacing a deterministic component that works
with a probabilistic one that can emit malformed Cypher. That is a different
project, and a worse design for this use case.

\begin{keybox}{{The one-sentence answer for the defence}}
``GraphRAFT applies when the language model generates the graph queries. In this
architecture traversal is deterministic and executed in code, so the trainable
component is passage \emph{selection} --- which is precisely what RAFT addresses.''
\end{keybox}

\subsection{Why traversal contributed nothing in this benchmark}

Of the 79 held-out questions where retrieval returned the answering passage,
\textbf{all 79 came from vector search and none from graph traversal}. Read
carelessly, that looks damning. It is not, and it must not be presented as a
finding.

\textbf{The benchmark cannot detect a traversal benefit, by construction.} Every
question was generated from \emph{one sentence in one Confluence document} --- a
single-hop lookup. Traversal exists to connect \emph{different} documents: which
ticket blocks another, which page is a child of which, which release caused a
regression. A single-hop question has nothing to traverse to, so the graph is
being measured on the one task where it cannot help.

\begin{warnbox}{{What is and is not safe to say about the graph}}
\textbf{Safe:} ``On single-hop documentation lookups, graph expansion added no
coverage while costing 2.4\,s per query. Its contribution to multi-hop questions
was not evaluated in this study.''

\textbf{Not safe:} ``The graph does not contribute.'' That claim requires a
multi-hop benchmark this work did not build.

\textbf{Worth noting:} traversal does \emph{not} depend on the Sentries API. The
graph already contains 428k JiraTicket and 47k ConfluencePage nodes with typed
relationships, so a multi-hop evaluation is possible today --- the API only
supplies \emph{live} ticket state.
\end{warnbox}

\section{Every step we took, and why it mattered}

\subsection{Step 1 --- Build questions with a known answer}

1{,}568 questions, each generated \emph{from} a specific sentence in a specific
document. This is the foundation of everything: because we know which sentence
answers each question, we can check whether that sentence was retrieved and
whether the answer used it. Without a known answer, none of the metrics are
computable.

Contexts came from your real retrieval path --- vector search, typed-relationship
traversal, cross-encoder rerank --- so the benchmark reflects the deployed system.

\subsection{Step 2 --- Split before building anything}

500 questions held out, keyed on question id, written to disk \emph{before} a single
training example was constructed.

\textbf{Why it matters:} if the split happened afterwards, a question could land in
both training and evaluation, and any improvement would be memorisation. Splitting
on \emph{question} rather than on \emph{answer} also matters --- we had two answers
per question from two engines, and splitting on answers would have put a question
in training and a near-copy in the holdout.

\subsection{Step 3 --- Recover the golden passage}

Retrieval returned the source document only 88\,\% of the time in the training
pool, so we fetched the golden passage back from Neo4j by title.

\textbf{Why it matters:} RAFT \emph{composes} its training context deliberately.
Relying on retrieval to supply the golden passage would have silently dropped the
majority of training examples. Evaluation was left untouched --- the holdout is
answered against whatever retrieval actually returns.

\subsection{Step 4 --- Generate reasoning chains with a stronger model}

993 chain-of-thought answers written by \texttt{gpt-oss:120b}, each validated:
any answer quoting text absent from the context was discarded (75 were, 7\,\%).

\textbf{Why it matters, and the trap avoided:} the teacher writes \emph{training
data only}. Evaluation is entirely deterministic --- validator plus refusal
classifier, no model judgement. If the same model both taught and graded, the
result would be circular and a jury would say so.

\subsection{Step 5 --- Compose the dataset}

914 training / 79 evaluation examples at $P = 0.80$, with the golden passage
inserted at a \textbf{random rank} among the distractors.

\textbf{Why the random rank matters:} if the golden passage were always first, the
model would learn ``the answer is at the top'' --- an artefact of our construction,
not a reading skill. Randomising forces it to actually locate the passage.

\subsection{Step 6 --- Measure the baseline before training}

500 held-out questions answered by the untouched model, scored on nine metrics
fixed in advance.

\textbf{Why in advance:} metrics chosen after seeing results are worthless. Ours
were written into the harness before the adapter existed, along with the numeric
thresholds that would decide whether to train at all.

\section{The metrics, explained}

\begin{table}[h]\centering\small
\begin{tabularx}{\textwidth}{l Y}
\toprule
\textbf{Metric} & \textbf{What it means and why it is there} \\
\midrule
\textbf{golden-fact recall} & Of the versions and quantities in the answering
sentence, how many appear in the model's answer. Measured only where retrieval
actually returned that sentence --- the model cannot use what it never saw.
\emph{This is the primary metric.} \\
\rowcolor{rowgray} \textbf{zero recall} & Share of answers containing \emph{none} of
those facts. The blunt version of the above: the model had the passage and used
nothing from it. \\
\textbf{recall by rank} & Recall broken down by where the answering passage sat in
the context. Separates what retrieval ordering causes (fixable in config) from what
the model causes (not). \\
\rowcolor{rowgray} \textbf{fabrication rate} & Validator v1.2: a version, config key,
quantity, page title or URL asserted but absent from the context. Deterministic. \\
\textbf{quote grounding} & Share of the answer's 8-word spans appearing verbatim in
the context. This measures RAFT's \emph{mechanism} directly, not just its outcome. \\
\rowcolor{rowgray} \textbf{refusal rate} & Guard metric. Must not rise. Adapter~1
``improved'' by declining more (75.4\,\% $\rightarrow$ 93.4\,\%); a metric that can be
gamed by silence is not a metric. \\
\textbf{bare vs helpful refusal} & Whether a refusal offers anything usable. \\
\rowcolor{rowgray} \textbf{evidence similarity} & Embedding similarity between answer
and answering sentence. Catches answers that cite the right version but describe it
wrongly. \\
\textbf{answer length} & Regression guard --- RAFT answers are longer by design;
tripling would be a problem, not a win. \\
\bottomrule
\end{tabularx}
\end{table}

Every one is deterministic and free. That is what makes a 500-question holdout
affordable, and a 3.5-point improvement detectable where a judge-scored holdout of
130 could only see 10 points.

\section{Baseline results, with presentation flags}

\begin{table}[h]\centering\small
\begin{tabularx}{\textwidth}{Y r l}
\toprule
\textbf{Result (500-question holdout)} & \textbf{Value} & \textbf{Flag} \\
\midrule
Fabrication rate & 2.6\,\% & \SHOW \\
\rowcolor{rowgray} Golden-fact recall when passage retrieved & 69.2\,\% & \SHOW \\
Perfect recall & 64.6\,\% & \SHOW \\
\rowcolor{rowgray} Declines when the passage is genuinely absent & 73.5\,\% & \SHOW \\
Zero recall at rank 1--2 & 28.4\,\% $\pm$ 10.8 & \CARE \\
\rowcolor{rowgray} Quote grounding & 2.19\,\% & \SHOW \\
Refusal rate & 66.4\,\% & \CARE \\
\rowcolor{rowgray} Bare refusals (of refusals) & 71.4\,\% & \CARE \\
Retrieval coverage & 14--16\,\% & \CARE \\
\rowcolor{rowgray} Primary metric sample size & $n=79$ of 500 & \HIDE{} \emph{as a headline} \\
Graph traversal supplied the answering passage & 0 of 79 & \HIDE \\
\bottomrule
\end{tabularx}
\end{table}

\subsection*{Why each flag}

\textbf{\SHOW{} --- Fabrication 2.6\,\%.} Low, deterministic, measured on 500
answers. This is a genuinely good result: your assistant rarely invents things.

\textbf{\SHOW{} --- Recall 69.2\,\%, perfect 64.6\,\%.} When the system works end to
end, it works well. Solid sample reasoning, clearly defined metric.

\textbf{\SHOW{} --- Declines 73.5\,\% when evidence is absent.} This is
\emph{correct behaviour} and worth presenting as such. Safety-first failure mode.

\textbf{\SHOW{} --- Quote grounding 2.19\,\%.} A clean baseline on the exact
mechanism RAFT trains. Whatever the adapter does to this number is meaningful.

\textbf{\CARE{} --- Zero recall 28.4\,\%.} The point estimate cleared the
pre-registered threshold, but the confidence interval runs 17.6--39.2\,\%. Present
it \emph{with} the interval. Never present it as ``28\,\%'' alone.

\textbf{\CARE{} --- Refusal 66.4\,\%.} True but easy to misread as the model being
unhelpful. It must be presented together with the coverage number, because most of
those refusals are correct given that the evidence was absent.

\textbf{\CARE{} --- Coverage 14--16\,\%.} Real and important, but measured on
\emph{templated} questions (``What does the X documentation state about Y''). Real
engineers do not phrase questions that way. Say ``on this benchmark'', never ``our
RAG retrieves badly''.

\textbf{\HIDE{} --- $n=79$ as a headline.} If you put the primary metric on a slide
without the sample size, the first question will expose it. Either show $n$ and the
interval, or lead with a different metric.

\textbf{\HIDE{} --- ``graph traversal contributed nothing''.} In all 79 cases the
answering passage came from vector search, none from traversal. \textbf{Do not
present this as a finding.} Every question in this benchmark is single-hop by
construction --- built from one sentence in one document --- which is exactly the
question type where traversal cannot help. The benchmark cannot evaluate the graph,
and the honest position is that its contribution here is untested, not zero.

\begin{warnbox}{{The retrieval sweep: what it does and does not license}}
We swept nine retrieval configurations. Coverage moved from 14.0\,\% to 14.8\,\% ---
essentially nothing --- across wider windows, more hops, larger expansion and no
reranker.

\textbf{Safe to say:} on this benchmark, retrieval coverage is not recoverable by
configuration; it points at chunking and embeddings, which is a re-index.

\textbf{Not safe to say:} that graph expansion is useless. It cost 2.4\,s per query
for no coverage gain \emph{on single-hop documentation questions}. Your Jira and
GitLab traffic --- where relationships between tickets are the whole point --- was
never tested, because the API was not available.
\end{warnbox}

\section{What happened, and what it taught}

Two adapters were trained on the first version of the RAFT set. \textbf{The gate
rejected both}, for opposite reasons, and the pair of failures is more instructive
than a single promotion would have been.

\begin{table}[h]\centering\small
\begin{tabularx}{\textwidth}{l Y}
\toprule
\textbf{Candidate} & \textbf{Gate verdict} \\
\midrule
Higher lr, full data & \textbf{REJECT} --- degenerate abstention: refusals rose past
the pre-registered limit \\
\rowcolor{rowgray} Lower lr, 73\,\% of data & \textbf{REJECT} --- fabrication rose:
refusals \emph{fell}, and the model answered where it should have declined \\
\bottomrule
\end{tabularx}
\end{table}

\begin{warnbox}{{The defect, and why nothing else caught it}}
Every golden-absent example shared \textbf{one hardcoded abstention string} --- 20\,\%
of the training set as a single identical target, and by far its largest duplicate.
The model did not learn to decline when evidence is missing; it learned the string,
and recited it verbatim on held-out questions whose evidence was present.

The loss curve looked healthy throughout. The training gates in place at the time
all passed. What caught it was the promotion gate, on the same
degenerate-abstention rule that had already rejected an earlier DPO adapter.

The identical recitation appeared on two independent stacks --- a rented T4 running
fp16, and the local server running 4-bit --- which is what establishes it as a
property of the data rather than of any run.
\end{warnbox}

This is the documented failure mode of refusal-aware instruction tuning: supervised
training on a fixed refusal pattern memorises the pattern rather than the decision
boundary, and over-refuses on answerable questions. The corrective is that an
abstention must be \emph{derived from its own context} --- naming the releases the
retrieved passages actually cover, and the one they do not --- so that no fixed
string exists to be memorised.

\begin{table}[h]\centering\small
\begin{tabularx}{\textwidth}{Y r r}
\toprule
\textbf{Training set} & \textbf{v1 (rejected)} & \textbf{v2 (current)} \\
\midrule
Examples & 914 & 912 \\
\rowcolor{rowgray} Unique targets & 735 / 914 & \textbf{912 / 912} \\
Most-repeated target & \textbf{180$\times$ (20\,\%)} & \textbf{1$\times$ (0.1\,\%)} \\
\rowcolor{rowgray} Quote-marker coverage & 72\,\% & 78\,\% \\
Spans marked verbatim that were not & 42 & \textbf{0} \\
\bottomrule
\end{tabularx}
\end{table}

Four pre-flight gates were added so the class of defect cannot recur silently: no
target may exceed 5\,\% of the set; no example may lose its answer to truncation;
training aborts when the loss stops being a finite number; and an adapter containing
NaN values or dead low-rank factors cannot be saved. Each exists because something
got past the checks that preceded it.

\begin{keybox}{{The sentence to say about this}}
The gate refused two adapters for measured reasons, and the second refusal is what
exposed a data defect that would otherwise have shipped. The deliverable is not one
adapter --- it is a loop that measures its own output, declines to deploy a
regression, and says why.
\end{keybox}

\section{Limitations --- state these before anyone asks}

\begin{enumerate}[leftmargin=1.4em,itemsep=3pt]
  \item \textbf{The primary metric rests on 79 questions.} Retrieval supplied the
        answering passage in only 16\,\% of cases, so the recall metric is only
        computable there. Confidence interval $\pm$10 points.
  \item \textbf{Questions are templated and machine-generated.} They are answerable
        by construction and phrased uniformly. Real user questions differ, and the
        coverage figure in particular may not transfer.
  \item \textbf{Single-hop only.} The benchmark cannot evaluate graph traversal,
        multi-document reasoning, or the live-API path.
  \item \textbf{The Sentries API was unavailable.} Everything here is Confluence
        documentation. Ticket-shaped questions were excluded because their
        refusals reflect a missing API, not model behaviour.
  \item \textbf{Training data is teacher-generated.} The reasoning chains were
        written by \texttt{gpt-oss}, not by humans or by the model itself. Quotes
        were validated against the context, but the reasoning was not human-checked.
  \item \textbf{Train/serve quantisation differs.} Training used 4-bit NF4;
        deployment through Ollama uses Q4\_K\_M. Measured earlier: fp16 fabricates
        \textasciitilde3 points less than Q4, so quantisation is not free.
  \item \textbf{One run, one seed.} No repeated runs, so run-to-run variance is
        unmeasured.
\end{enumerate}

\section{Defence preparation}

\begin{table}[h]\centering\small
\begin{tabularx}{\textwidth}{Y Y}
\toprule
\textbf{Likely question} & \textbf{Answer} \\
\midrule
``Why not just fix the retrieval?'' & We tried. Nine configurations moved coverage
by 0.8 points. And the defect we trained on occurs when the passage is already
ranked first --- ordering cannot explain it. \\
\rowcolor{rowgray} ``Why not a filter for fabricated URLs?'' & We should, and it is
cheaper. It fixes fabricated citations. It cannot make a model read the paragraph
it was given, which is what the 28\,\% is. \\
``Isn't 79 questions too few?'' & For the primary metric, yes, and the interval is
stated. The mechanism metric (quote grounding) is measured on all 500. \\
\rowcolor{rowgray} ``Did the graph help?'' & Untested here. This benchmark is
single-hop by construction; traversal is designed for multi-hop ticket
relationships, which the missing API prevented us from evaluating. \\
``Why RAFT rather than DPO?'' & DPO was run first and the gate rejected it. The
pairs were 93\,\% synthetic corruptions, and the model learned to decline rather
than to ground. RAFT targets document selection with real reasoning chains. \\
\rowcolor{rowgray} ``What if the gate rejects this one too?'' & Then the loop worked.
A gate that never rejects is decoration. The first rejection was diagnosed --- refusal
contamination in the pairs --- and the pipeline corrected. \\
\bottomrule
\end{tabularx}
\end{table}

\section*{References}
\small
\begin{itemize}[leftmargin=1.2em,itemsep=1pt]
  \item Zhang, Patil et al. \textit{RAFT: Adapting Language Model to Domain Specific RAG}, 2024. \url{https://arxiv.org/abs/2403.10131}
  \item \textit{Divide-Then-Align: Honest Alignment based on the Knowledge Boundary of RAG}, ACL 2025. \url{https://arxiv.org/pdf/2505.20871}
  \item Joren et al. \textit{Sufficient Context: A New Lens on Retrieval Augmented Generation}, 2024. \url{https://arxiv.org/pdf/2411.06037}
  \item Hu et al. \textit{LoRA: Low-Rank Adaptation of Large Language Models}, 2021.
  \item Dettmers et al. \textit{QLoRA: Efficient Finetuning of Quantized LLMs}, 2023.
  \item Rafailov et al. \textit{Direct Preference Optimization}, 2023.
  \item Meta AI. \textit{RAFT: Sailing Llama towards better domain-specific RAG}, 2024.
\end{itemize}

\end{document}
